\subsection{The full \texorpdfstring{\(S^{12}\to S^6\)}{S12-to-S6}
carrier is null-homotopic}
\label{cusp-s12-carrier-homotopy}

This subsection determines the based homotopy class of the full two-copy
sphere carrier.  It is important that the calculation concerns the full
quotient map, not merely the inclusion of one of its moving four-spheres.

Recall the exact maps from
\cref{cusp-cone-S12-route,cusp-cone-angle-forgotten-sphere-carriers}.  The
cusp quotient is
\[
 q_2:(X,x_0)\longrightarrow
 \mathcal A=\operatorname{Susp}(S^1\times\overline T^2),
\]
and, for the coordinate one-skeleton
\[
 K^{(1)}=
 ((\mathbb R/\mathbb Z)\times\{0\})
 \cup(\{0\}\times(\mathbb R/\mathbb Z))\subset\overline T^2,
\]
the degree-one quotient
\(c:\overline T^2\to\overline T^2/K^{(1)}\cong S^2\) defines
\[
 \widehat p_2([r,e^{i\phi},u])=[r,c(u)]
 \in\Sigma_r S^2\cong S^3.
\]
Thus
\[
 g:=\widehat q_3=\widehat p_2\circ q_2:(X,x_0)\longrightarrow(S^3,*)
\]
is the specified continuous surjective based quotient.  After the fixed
based identification \(X\cong_{\mathrm{diff}}S^6\), its smash square is the
specified continuous surjective based quotient
\[
 g^{\wedge2}:S^{12}=S^6\wedge S^6
 \longrightarrow S^3\wedge S^3=S^6.
\]

\begin{theorem}[The full external smash pairing is zero]
\label{cusp-s12-full-carrier-null}
For every pair of based maps \(x,y:S^6\to S^3\),
\[
 [x\wedge y]=0\qquad\text{in}\qquad\pi_{12}(S^6).
\]
In particular,
\(\widehat q_3^{\wedge2}:S^{12}\twoheadrightarrow S^6\) is simultaneously
surjective as a point-set quotient and based null-homotopic.
\end{theorem}

\begin{proof}
Write \(E\) for reduced suspension, using the ordered smash coordinates on
spheres.  The external smash product has the exact factorization
\begin{equation}
 x\wedge y
 =
 (x\wedge\operatorname{id}_{S^3})
 \circ
 (\operatorname{id}_{S^6}\wedge y):
 S^6\wedge S^6\longrightarrow S^6\wedge S^3
 \longrightarrow S^3\wedge S^3.
 \label{cusp-s12-smash-factorization}
\end{equation}
Under the displayed ordered sphere identifications, the two factors, in
order of application, represent \(E^6y\in\pi_{12}(S^9)\) and
\(E^3x\in\pi_9(S^6)\).  Equivalently, Toda's
Barratt--Hilton reduced-join formula gives, with
\(p=q=k=h=3\),
\begin{equation}
 [x\wedge y]=(E^3[x])\circ(E^6[y]).
 \label{cusp-s12-toda-smash-formula}
\end{equation}
Here Toda's sign is \((-1)^{(p+k)h}=(-1)^{18}=+1\).  Reversing the
two displayed factors introduces the separate twist sign
\((-1)^{ph}=-1\); no such reversal is used here.

The three exact groups needed below are
\begin{equation}
 \pi_6(S^3)\cong\mathbb Z/12,
 \qquad
 \pi_{12}(S^9)\cong\mathbb Z/24,
 \qquad
 \pi_{12}(S^6)\cong\mathbb Z/2.
 \label{cusp-s12-three-stem-groups}
\end{equation}
The first and third are the integral low unstable groups displayed in
Hatcher's table extracted from Toda; the middle group is in the stable
three-stem
\cite[opening table, p.~339]{Hatcher2002};
\cite[Table~1, p.~113]{IsaksenWangXu2023}.

Put \(\beta=E^6[y]\in\pi_{12}(S^9)\).  Since
\(12[y]=0\), functoriality of suspension gives \(12\beta=0\).  In the
cyclic group \(\mathbb Z/24\), the elements annihilated by \(12\) are
exactly the even elements.  Hence there is
\(\gamma\in\pi_{12}(S^9)\) such that \(\beta=2\gamma\).
Postcomposition with \(E^3[x]:S^9\to S^6\) induces a group homomorphism
\[
 (E^3[x])_*:\pi_{12}(S^9)\longrightarrow\pi_{12}(S^6).
\]
The target has exponent two, so \cref{cusp-s12-toda-smash-formula} gives
\[
 [x\wedge y]
 =(E^3[x])_*(\beta)
 =2(E^3[x])_*(\gamma)
 =0.
\]
This proves the vanishing of the entire external smash pairing and hence
the theorem.
\end{proof}

\begin{corollary}[All higher smash powers]
\label{cusp-s12-all-higher-smashes-null}
For every based \(g:S^6\to S^3\) and every \(r\geq2\),
\[
 g^{\wedge r}:S^{6r}\longrightarrow S^{3r}
\]
is based null-homotopic.
\end{corollary}

\begin{proof}
With the fixed associativity homeomorphisms for smash products,
\[
 g^{\wedge r}=g^{\wedge2}\wedge g^{\wedge(r-2)}.
\]
Smashing a based null-homotopy with a based map gives a based
null-homotopy.
\end{proof}

\begin{remark}[The one-copy invariant and its computed value]
\label{cusp-s12-one-copy-invariant}
The proof does not need to determine the one-copy class
\([g]\in\pi_6(S^3)\cong\mathbb Z/12\).  Toda's suspension calculation
identifies \(\pi_6(S^3)\) with the even subgroup of the stable three-stem;
Hatcher gives an integral Serre-spectral-sequence derivation of the same
injectivity and identifies twice the stable Hopf--\(\nu\) generator with a
generator of the suspended \(\mathbb Z/12\)-summand
\cite[Theorem~5.41(b), p.~573, and pp.~577--578]{HatcherSpectralSequences}.
Thus, after choosing a stable generator \(\nu_9\), there is a unique
\(a\in\mathbb Z/12\) such that
\begin{equation}
 [E^6g]=2a\nu_9\in\pi_{12}(S^9)\cong\mathbb Z/24.
 \label{cusp-s12-sixfold-suspension-class}
\end{equation}
Because \(a\mapsto2a\) embeds \(\mathbb Z/12\) as the even subgroup of
\(\mathbb Z/24\), \(E^6g\) is null-homotopic if and only if \(g\) is
null-homotopic.  Thus the one-copy class and its sixfold suspension are the
same integer modulo \(12\). It is now zero by Theorem~\ref{cusp-null:thm}.

The exact geometric invariant which computes \(a\) is also identifiable.
Choose \(0<r_0<\epsilon\) and
\(u_0\in\overline T^2\setminus K^{(1)}\), and let
\(y_0=[r_0,c(u_0)]\) lie in the open three-cell of
\(\Sigma_rS^2\).  From the literal formula for \(q_2\),
\begin{equation}
 g^{-1}(y_0)
 =\{x\in f^{-1}(|t|=r_0):\overline\pi_f(x)=u_0\}.
 \label{cusp-s12-regular-inverse-image}
\end{equation}
The fibre of \(\overline\pi\) is the two-torus defined by
\(K_\Lambda=\mathbb Z\langle\widehat w,\widehat\delta\rangle\), and the
cusp monodromy fixes both \(\widehat w\) and \(\widehat\delta\).  Therefore
\eqref{cusp-s12-regular-inverse-image} is the mapping torus of an affine translation of the kernel torus,
which is diffeomorphic to \(S^1\times T^2\cong T^3\): if the translation is
\(k\mapsto k+v\), choose a real lift \(\widetilde v\); the map
\([\theta,k]\mapsto(e^{2\pi i\theta},k+\theta\widetilde v)\) respects
\((1,k)\sim(0,k+v)\) and has the inverse obtained by subtracting
\(\theta\widetilde v\). After choosing a smooth chart in the open
target cell, the differential of \(g\) supplies its normal three-framing.
Pontryagin--Thom identifies the stabilized class of \(g\) with the framed
bordism class of this specifically framed \(T^3\).  Since stabilization on
\(\pi_6(S^3)\) is injective, computing that framing class computes \(a\)
exactly.  The unframed diffeomorphism type \(T^3\), point-set surjectivity,
and the zero ordinary homology map do not compute \(a\).
\end{remark}

\begin{remark}[Scope]
\label{cusp-s12-nullity-scope}
The theorem changes no point-set or metric statement.  A surjective quotient
map may be null-homotopic, and the embedded moving four-spheres and their
intrinsic pointed \(\mathbb R^4\)-limits remain exactly as constructed.  The
theorem concerns the sphere-valued postcomposition
\(\widehat q_3^{\wedge2}\); it does not imply null-homotopy of the earlier
conical-target map
\(q_2^{\wedge2}:S^{12}\to\mathcal A^{\wedge2}\).
\end{remark}

\subsection{The original cusp coefficient: an explicit framed nullbordism}


Let \(X\) be the smooth compact source of the original construction,
\(f:X\to B\) its map, and \(t\) the specified cusp coordinate.
Put \(W=f^{-1}(0)\), choose the original \(x_0\in W\), and retain the
ordered quotient lattice basis induced by
\((\widehat\gamma,\widehat u)\), with kernel generated by
\((\widehat w,\widehat\delta)\). Thus the quotient torus is
\(\overline T^2=\mathbb R^2/\mathbb Z^2\). The original periods are
\begin{equation}
 Z(s)=sB_0+C(t),\qquad e^{2\pi i s}=t,\qquad
 B_0=\begin{pmatrix}0&1\\-1&0\end{pmatrix},\qquad
 C(t)=\begin{pmatrix}6\mu(t)&h(t)\\b(t)-h(t)&\mu(t)\end{pmatrix}.
 \label{cusp-null:eq:periods}
\end{equation}
Every entry of \(C\) is holomorphic at \(t=0\).
The dense torus of the original toric cusp has coordinates
\((x_1,x_2,t)\), with
\(x=\exp(2\pi i\zeta)\). The residual lattice action is exactly
\begin{equation}
 \Psi_\lambda(x,t)=
 \bigl(t^{B_0\lambda}\exp(2\pi iC(t)\lambda)x,t\bigr),
 \qquad\lambda\in\mathbb Z^2.
 \label{cusp-null:eq:action}
\end{equation}
It extends freely and properly discontinuously across the original
unimodular toric filling. In particular the quotient projection is a local
biholomorphism at every toric fixed point. These are the reconstructed
data and proved local properties of the original source
(the reconstructed original manuscript, \S4, source pp.~24--29; its complete cusp data are displayed here).

Let
\[
 K^{(1)}=(\mathbb R/\mathbb Z\times\{0\})\cup(\{0\}\times\mathbb R/\mathbb Z),\qquad
 c:\overline T^2\longrightarrow \overline T^2/K^{(1)}\cong S^2
\]
be the specified degree-one cell quotient. Its restriction to the open
two-cell is taken with the smooth coordinates used in the original
regular-fibre description. Both radial endpoints and the complete
basepoint cylinder are collapsed in
\[
 \Sigma_r(S^2,*)=
 \frac{[0,\epsilon]\times S^2}
 {(\{0,\epsilon\}\times S^2)\cup([0,\epsilon]\times\{*\})}
 \cong S^3.
\]
The original based map is
\begin{equation}
 g(x)=[|t|,c(\overline\pi_f(x))]\quad(0<|t|<\epsilon),\qquad
 g(W)=g(X\setminus\operatorname{int}N_0)=*.
 \label{cusp-null:eq:g}
\end{equation}

\begin{theorem}[The original one-copy cusp map is null-homotopic]
\label{cusp-null:thm}
The based map \(g:(X,x_0)\to(S^3,*)\) in \eqref{cusp-null:eq:g} is based
null-homotopic. More precisely, a regular inverse image with its actual
ordered normal framing induced by \((r,c_1,c_2)\) is framed nullbordant
inside one original cusp chart times \([0,1]\), disjoint from
\(\{x_0\}\times[0,1]\).
\end{theorem}

\begin{proof}
The proof constructs the complete inverse image, compares its actual
framing with an explicitly extendible framing, and gives the extension.

\paragraph{1. The exact complete inverse image.}
Retain the lift
\[
 u^0=\begin{pmatrix}-1/4\\1/4\end{pmatrix},\qquad
 \overline u^0=(3/4,1/4)\pmod{\mathbb Z^2},\qquad
 B_0u^0=\begin{pmatrix}1/4\\1/4\end{pmatrix}.
\]
Thus \(\overline u^0\notin K^{(1)}\), and \(B_0u^0\) lies in the
interior of the original triangle \(\operatorname{conv}(0,e_1,e_2)\).
Use its original unimodular chart
\begin{equation}
 z_0=t/(x_1x_2),\qquad z_1=x_1,\qquad z_2=x_2,
 \qquad t=z_0z_1z_2.
 \label{cusp-null:eq:chart}
\end{equation}
Set
\[
 \ell=\log|t|<0,\qquad R(t)=-2\pi\operatorname{Im}C(t),\qquad
 A(t)=\ell B_0+R(t).
\]
For a representative \(\zeta=Z(s)u+k\), where \(u,k\in\mathbb R^2\),
\(\operatorname{Im}s=-\ell/(2\pi)\) gives the exact identity
\begin{equation}
 \begin{pmatrix}\log|x_1|\\\log|x_2|\end{pmatrix}
 =A(t)u.
 \label{cusp-null:eq:moduli}
\end{equation}
The quotient value is \(u\bmod\mathbb Z^2\).
Since \(B_0\) is orthogonal, for all \(v\in\mathbb R^2\),
\begin{equation}
 \|A(t)v\|\geq (|\ell|-\|R(t)\|)\|v\|.
 \label{cusp-null:eq:invertibility}
\end{equation}
The boundedness of \(R\) therefore makes \(A(t)\) invertible on a
sufficiently small punctured disc.

For any point of the fibre quotient over \(\overline u^0\), subtracting
the appropriate original period, equivalently applying
\(\Psi_{-\lambda}\), gives a representative with the exact lift
\(u=u^0\). This representative is unique: another residual translate
has logarithmic moduli differing by \(A(t)\lambda\), which vanishes
only for \(\lambda=0\). Conversely every dense-torus point satisfying
\(\log|x|=A(t)u^0\) has this quotient value, since the remaining real
parts of \(\zeta\) are exactly the kernel coordinates \(k\).
For each fixed \(t\), varying \(k\bmod\mathbb Z^2\) gives all independent
phases of \(x_1,x_2\).

The following change of coordinates uses the full original period units:
\begin{align}
 h_j^0(t)&=\exp\bigl(-2\pi i(C(t)u^0)_j\bigr),\qquad j=1,2,\\
 \mathcal F(z_0,z_1,z_2)
 &=(w_0,w_1,w_2)
 =\left(\frac{z_0}{h_1^0(t)h_2^0(t)},\,
 h_1^0(t)z_1,\,h_2^0(t)z_2\right).
 \label{cusp-null:eq:straightening}
\end{align}
The functions \(h_j^0\) are nowhere-zero holomorphic units, including
at \(t=0\). The product \(w_0w_1w_2\) is exactly \(t\).
The inverse of \(\mathcal F\) is
\[
 (w_0,w_1,w_2)\longmapsto
 \left(w_0h_1^0(t)h_2^0(t),\,
 \frac{w_1}{h_1^0(t)},\,\frac{w_2}{h_2^0(t)}\right),
 \qquad t=w_0w_1w_2.
\]
Thus this is a biholomorphism of the original chart restricted to the
cusp disc, not a replacement of \(C(t)\).

For \(r_0\in(0,\epsilon)\), the entire inverse image of
\(y_0=[r_0,c(\overline u^0)]\) becomes exactly
\begin{equation}
 T_{r_0}=
 \{|w_0|=\rho_0,\ |w_1|=\rho_1,\ |w_2|=\rho_2\},\qquad
 (\rho_0,\rho_1,\rho_2)=
 (r_0^{1/2},r_0^{1/4},r_0^{1/4}).
 \label{cusp-null:eq:torus}
\end{equation}
Indeed, \(\log|h^0(t)|=2\pi\operatorname{Im}C(t)u^0=-R(t)u^0\),
so \eqref{cusp-null:eq:moduli} gives \(\log|w_j|=\ell(B_0u^0)_j\) for
\(j=1,2\); the product identity gives the exponent \(1/2\) for
\(w_0\). Conversely these moduli recover \(u=u^0\) by invertibility
of \(A(t)\). The angular map is explicitly
\[
 (\arg t,\arg w_1,\arg w_2)\longleftrightarrow
 (\arg w_0,\arg w_1,\arg w_2),\qquad
 \arg w_0=\arg t-\arg w_1-\arg w_2\pmod{2\pi},
\]
with inverse \(\arg t=\arg w_0+\arg w_1+\arg w_2\).
Its integral matrix has determinant one. This includes the complete
angular mapping torus and all kernel phases, with their original signs.

Let \(o\) be the origin of \eqref{cusp-null:eq:chart}. The original free properly
discontinuous action supplies a neighborhood \(V\) of \(o\) on which
the quotient projection is injective. To see the local injectivity
directly, first choose a relatively compact neighborhood; proper
discontinuity leaves only finitely many translates that can intersect it.
None of those nonidentity translates fixes \(o\). Shrinking finitely
many times separates the neighborhood from all of them. The quotient
map is a local biholomorphism, so its restriction to \(V\) is a chart
of the original source.
Choose a polydisc about \(0\) whose closure lies in \(\mathcal F(V)\).
For sufficiently small \(r_0\), both \(T_{r_0}\) and
\begin{equation}
 \{|w_0|\leq\rho_0,\ |w_1|=\rho_1,\ |w_2|=\rho_2\}
 \label{cusp-null:eq:spatial-bound}
\end{equation}
lie in this polydisc. This proves that the whole inverse image embeds
in one chart and supplies space for the bounding manifold below.

\paragraph{2. The actual normal framing, including its angle dependence.}
Put \(L_j=\log|w_j|\), and use the unique real lift of the quotient
coordinate near \(u^0\). From the exact change of coordinates,
\begin{equation}
 \begin{pmatrix}L_1\\L_2\end{pmatrix}
 =\ell B_0u^0+A(t)(u-u^0),\qquad
 L_0=\ell-L_1-L_2.
 \label{cusp-null:eq:nearby-moduli}
\end{equation}
Differentiating on \(u=u^0\) annihilates the term
\((dA(t))(u-u^0)\). Hence no derivative of \(C(t)\) has been omitted:
the exact differential there is
\begin{equation}
 \begin{pmatrix}dL_0\\dL_1\\dL_2\end{pmatrix}
 =J(t)\begin{pmatrix}d\ell\\du_1\\du_2\end{pmatrix},\quad
 J(t)=
 \begin{pmatrix}
 1/2&-A_{11}-A_{21}&-A_{12}-A_{22}\\
 1/4&A_{11}&A_{12}\\
 1/4&A_{21}&A_{22}
 \end{pmatrix}.
 \label{cusp-null:eq:J}
\end{equation}
Adding the second and third rows to the first gives \((1,0,0)\),
and expansion in that row proves \(\det J(t)=\det A(t)\).

Let \(D=Dc_{\overline u^0}\) be the actual two-by-two derivative in
the ordered open-cell target coordinates \((c_1,c_2)\). It is a fixed
invertible matrix; all its entries and signs are retained. Let
\[
 H=\begin{pmatrix}r_0&0&0\\0&D_{11}&D_{12}\\0&D_{21}&D_{22}\end{pmatrix}.
\]
Let \(n_j\) be the real radial vector \(w_j\) in the \(j\)-th complex
coordinate, with other components zero. Then \(dL_i(n_j)=\delta_{ij}\).
Writing \(N=(n_0,n_1,n_2)\) as the matrix with these vector columns, the actual
\(dg\)-dual normal framing is, in the normal quotient bundle,
\begin{equation}
 \mathfrak n(t)=N J(t)H^{-1}.
 \label{cusp-null:eq:actual-frame}
\end{equation}
Indeed, \(dg(N)=HJ(t)^{-1}\), so multiplication by \(J(t)H^{-1}\)
gives the identity. In the Euclidean metric of this chart the radial
vectors are perpendicular to the product torus; a different source
metric changes their representatives by the canonical projection to its
normal bundle and preserves this normal-quotient framing.

Put \(\ell_0=\log r_0\). Decrease \(r_0\) further, if necessary, so
that \(|\ell_0|>2\sup_{|t|\leq r_0}\|R(t)\|\). For
\(0\leq\sigma\leq1\), define only a homotopy of frames by
\[
 A_\sigma(t)=\ell_0B_0+\sigma R(t),
\]
and let \(J_\sigma(t)\) be \eqref{cusp-null:eq:J} with \(A_\sigma\) in place
of \(A\). For every \(v\in\mathbb R^2\),
\[
 \|A_\sigma(t)v\|\geq
 (|\ell_0|-\|R(t)\|)\|v\|>0\quad(v\ne0).
\]
Thus all matrices are invertible. At \(\sigma=0\),
\begin{equation}
 J_0=
 \begin{pmatrix}
 1/2&\ell_0&-\ell_0\\
 1/4&0&\ell_0\\
 1/4&-\ell_0&0
 \end{pmatrix},\qquad \det J_0=\ell_0^2>0.
 \label{cusp-null:eq:J0}
\end{equation}
Continuity also proves \(\det J_\sigma(t)>0\) for every \(\sigma,t\).
The homotopy \(N J_\sigma(t)H^{-1}\), run from \(\sigma=1\) to
\(\sigma=0\), compares the actual frame with the exact constant-matrix
frame \(N P\), where
\begin{equation}
 P=J_0H^{-1}.
 \label{cusp-null:eq:P}
\end{equation}
The original embedded torus and original functions \(C(t)\) have not
changed. A framing homotopy is a framed cylinder: choose a smooth
parameter function constant near its endpoints and use the displayed
frames on \(T_{r_0}\times[0,1]\).

\paragraph{3. A collared framed nullbordism with the exact constant matrix.}
Use a new time coordinate \(v\in[0,1]\). Define
\begin{equation}
 q(v)=
 \begin{cases}
 0,&0\leq v\leq1/4,\\
 \rho_0^2\exp\bigl(4-1/(v-1/4)\bigr),&1/4<v\leq1.
 \end{cases}
 \label{cusp-null:eq:q}
\end{equation}
This function is smooth: all right derivatives of
\(\exp(-1/(v-1/4))\) at \(1/4\) vanish. It is strictly increasing
for \(v>1/4\), satisfies \(q(1/2)=\rho_0^2\), and
\(q'(1/2)=16\rho_0^2>0\).
In the chart times \([0,1]\), put
\begin{equation}
 \mathcal W=
 \{|w_0|^2+q(v)=\rho_0^2,\quad
       |w_1|^2=\rho_1^2,\quad|w_2|^2=\rho_2^2\}.
 \label{cusp-null:eq:bordism}
\end{equation}
No point has \(v>1/2\). The three defining differentials are independent:
the latter two have nonzero gradients in their separate complex factors;
the first has nonzero \(w_0\)-gradient whenever \(w_0\ne0\), and at
\(w_0=0\) it has the nonzero \(v\)-derivative \(q'(1/2)\).
Thus \(\mathcal W\) is a smooth compact four-manifold. Its only
boundary is \(T_{r_0}\times\{0\}\), and on \(0\leq v\leq1/4\)
it is the exact product collar \(T_{r_0}\times[0,1/4]\).
The north locus \(w_0=0,v=1/2\) is an interior cap: the implicit
function theorem expresses \(v\) smoothly in terms of
\((\operatorname{Re}w_0,\operatorname{Im}w_0)\) there. The first
factor is a disk obtained by rotating its radial profile with this smooth
cap and one boundary circle. Therefore
\(\mathcal W\cong D^2\times T^2\).
Its spatial coordinates lie in \eqref{cusp-null:eq:spatial-bound}, so the complete
manifold lies in the injective original chart.

In the Euclidean chart product metric define
\[
 V_0=(w_0,0,0,q'(v)/2),\qquad
 V_1=(0,w_1,0,0),\qquad V_2=(0,0,w_2,0).
\]
These are one half of the gradients of the three defining equations
in \eqref{cusp-null:eq:bordism}. They are mutually orthogonal and never zero.
For example, if \(F_0=|w_0|^2+q(v)-\rho_0^2\), then
\[
 dF_0(V_0)=2|w_0|^2+\frac12q'(v)^2>0
\]
at every point of \(\mathcal W\). On the entire collar \(q'=0\), so
\((V_0,V_1,V_2)=(N,0)\) exactly. Consequently
\begin{equation}
 (V_0,V_1,V_2)P
 \label{cusp-null:eq:extended-frame}
\end{equation}
is a framing of the normal three-plane bundle of \(\mathcal W\),
and its restriction to the collar is exactly the framing \(NP\) from
\eqref{cusp-null:eq:P}, with every original target-coordinate sign retained.
Concatenate the framing cylinder of the preceding paragraph with this
nullbordism, rescaling the two interval lengths. Their framings agree
exactly on the common product collar, so the concatenation is a smooth
framed nullbordism of the actual inverse image.

\paragraph{4. The based map represented by this framed inverse image.}
Choose the injective chart to exclude \(x_0\) if its center is not
\(x_0\). If its center is \(x_0\), the constructed nullbordism still
excludes \(x_0\), since \(|w_1|=\rho_1>0\) and
\(|w_2|=\rho_2>0\) everywhere on it. Compactness therefore gives
tubular neighborhoods for the framed cylinder and cap disjoint from
\(\{x_0\}\times[0,1]\).

The original continuous quotient map is smooth and a submersion near
its entire inverse image of \(y_0\). First make it constant in a
neighborhood of \(x_0\), by a based homotopy that changes nothing near
that inverse image. Explicitly, take a target chart
\(\kappa:O\to\mathbb R^3\) with \(\kappa(*)=0\), whose image is a ball,
and \(y_0\notin O\). By continuity choose a neighborhood \(U\) of
\(x_0\) with \(g(U)\subset O\), and a smooth bump function
\(\beta\) supported in \(U\) and equal to one near \(x_0\).
For \(x\in U\), replace \(g(x)\) during this homotopy by
\(\kappa^{-1}((1-a\beta(x))\kappa(g(x)))\), \(0\leq a\leq1\),
and leave \(g\) fixed off \(U\). Convexity of the target chart image
makes the formula well-defined, the vanishing of \(\beta\) near the
boundary of its support proves continuity of the pasted formula, and
\(\kappa(g(x_0))=0\) makes the homotopy based. It never introduces
an inverse image of \(y_0\). Its endpoint is smooth near \(x_0\) and
agrees with \(g\) near the entire old inverse image.

Denote the endpoint of the preceding homotopy by \(g_b\). Choose one
closed tubular neighborhood \(U_1\) of the inverse image, contained in
the smooth submersion region and with the entire inverse image contained
in \(\operatorname{int}U_1\). The compact set
\(X\setminus\operatorname{int}U_1\) has image under \(g_b\) missing
\(y_0\). In a fixed target metric put
\[
 d_0=\min_{x\in X\setminus\operatorname{int}U_1}
       d(g_b(x),y_0)>0.
\]
Relative smooth approximation, fixed on the whole closed set \(U_1\)
and on a neighborhood of \(x_0\), gives a smooth based representative
\(g_{\mathrm{sm}}\) whose error on
\(X\setminus\operatorname{int}U_1\) is less than \(d_0/2\).
It therefore has the same inverse image and differential on \(U_1\),
and \(d(g_{\mathrm{sm}}(x),y_0)>d_0/2\) on its complement. Thus there
is no uncontrolled intervening region and no additional inverse image.
Make the approximation also closer than a fixed target injectivity
radius; its short geodesics give the based homotopy.

The relative Pontryagin construction now applies to the explicit
framed nullbordism. It produces a based homotopy to a map missing
\(y_0\); stereographic contraction of \(S^3\setminus\{y_0\}\cong\mathbb R^3\)
to the image of \(*\) is based and contracts that last map.
For completeness, the dependence on the framed inverse image in this
argument is precisely the local product-neighborhood construction and
the interpolation lemma of \cite[\S7, Theorem B and Lemma 4,
pp.~43, 48--50]{CuspNullMilnor1965}. Their relative versions fix the source
basepoint because all tubular modifications are away from it and
stereographic interpolation between equal values keeps that value
fixed. A framed cylinder realizes a framing homotopy; the empty
Pontryagin manifold realizes a map missing the chosen regular value.
This proves the claimed based null-homotopy of the exact map \(g\).
\end{proof}

\begin{corollary}
Under the original fixed based identification \(X\cong S^6\), the
one-copy coefficient is
\[
 [g]=0\in\pi_6(S^3)\cong\mathbb Z/12,\qquad a=0\pmod{12},\qquad
 [E^6g]=0\in\pi_{12}(S^9)\cong\mathbb Z/24.
\]
The previously defined homomorphism
\(\Theta_g:\mathbb Z/12\to\pi_6(S^3)\), \(\Theta_g([n])=n[g]\), and
its composite with the marked peripheral character are zero.
\end{corollary}
\begin{proof}
The theorem states that the actual based map is null-homotopic.
Precomposition with the fixed based identification preserves a based
null-homotopy. Suspension preserves it as well. Thus the unique
coefficient of \([g]\) in any generator of \(\mathbb Z/12\) is zero.
For every \([n]\), \(\Theta_g([n])=n\cdot0=0\), and every
composite with this homomorphism is zero.
\end{proof}

\paragraph{Exact scope.}
The proof computes this sphere-valued postcomposition. It does not
compute the homotopy class of the earlier conical-target map
\(q_2:X\to\operatorname{Susp}(S^1\times\overline T^2)\), and it
does not assert a null-homotopy for that different target. Point-set
surjectivity of \(g\) and all of its stated metric restrictions remain
properties of the original map: a homotopy-class computation does not
alter the defining formula \eqref{cusp-null:eq:g}.
The retained-angle companion is proved separately in
Theorem~\ref{h24:ret:main}: $p_3q_2:X\to S^4$ represents
$\eta_4\eta_5\ne0$, so $q_2$ is nonnull. Its Hopf smash
represents $12\nu_6\ne0$ in $\pi_9(S^6)$. These positive
computations use the actual four-framed inverse torus of that different
postcomposition; they do not reverse or contradict the preceding
null-homotopy through $\widehat p_2$.


\begin{thebibliography}{9}
\bibitem{CuspNullMilnor1965}
J.~W. Milnor, \emph{Topology from the Differentiable Viewpoint},
University Press of Virginia, 1965, \S7, Theorem B and Lemma 4,
pp.~42--50. \url{https://www.maths.ed.ac.uk/~aar/papers/milnortop.pdf}.
\end{thebibliography}
